% Adaptacja merytoryczna AM I-09; nowe etykiety dla twierdzeń wektorowych.
\chapter{Szeregi w przestrzeniach Banacha}
\label{ma2:ch:banach-series}

W tym rozdziale $E$ jest przestrzenią unormowaną nad
$\K\in\{\R,\C\}$. Zupełność będziemy zakładać wyraźnie tam,
gdzie jest potrzebna. Wykorzystujemy rachunek szeregów liczbowych
z tomu I (rozdział 09) i wyniki rozdziału \ref{ma2:ch:normed}.
W szczególności zbieżność będzie rozumiana w normie, a nie
względem poszczególnych współrzędnych.

\section{Sumy częściowe i warunek Cauchy'ego}
\begin{definicja}[Szereg wektorowy]
\label{ma2:def:vector-series}
Dla ciągu $(x_n)_{n\ge1}$ w $E$ określamy sumy częściowe
$S_0=0$ i $S_N=\sum_{n=1}^N x_n$ dla $N\ge1$.
Szereg $\sum_{n=1}^\infty x_n$ nazywamy zbieżnym w $E$,
jeżeli istnieje $S\in E$ takie, że $\|S_N-S\|\to0$.
Piszemy wówczas $\sum_{n=1}^\infty x_n=S$.
W przeciwnym razie szereg jest rozbieżny.
\end{definicja}
Suma jest jednoznaczna: gdyby $\|S_N-S\|\to0$ i
$\|S_N-T\|\to0$, nierówność
$\|S-T\|\le\|S-S_N\|+\|S_N-T\|$ dawałaby $\|S-T\|=0$,
czyli $S=T$. Dla szeregu od indeksu $m\ge0$ stosujemy
analogiczną definicję, a sumę pustą przyjmujemy równą $0$.

\begin{twierdzenie}[Warunek konieczny i kryterium Cauchy'ego]
\label{ma2:thm:series-cauchy}
Jeżeli szereg $\sum_{n\ge1}x_n$ jest zbieżny w $E$, to
$x_n\to0$ i dla każdego $\varepsilon>0$ istnieje $N\ge1$
takie, że
\[
 \left\|\sum_{n=p}^{q}x_n\right\|<\varepsilon
 \quad\text{dla wszystkich całkowitych }q\ge p\ge N.
\]
Jeżeli $E$ jest przestrzenią Banacha, powyższy warunek
dotyczący sum skończonych jest również wystarczający.
\end{twierdzenie}
\begin{proof}
Niech $S_N\to S$. Dla $n\ge1$ mamy $x_n=S_n-S_{n-1}$,
więc
$\|x_n\|\le\|S_n-S\|+\|S_{n-1}-S\|\to0$.
Dla $\varepsilon>0$ istnieje $M$ takie, że
$\|S_j-S\|<\varepsilon/2$ dla $j\ge M$.
Jeśli $q\ge p\ge M+1$, to
\[
 \left\|\sum_{n=p}^q x_n\right\|=\|S_q-S_{p-1}\|
 \le\|S_q-S\|+\|S_{p-1}-S\|<\varepsilon.
\]
Odwrotnie, niech warunek sum skończonych zachodzi.
Dla $\varepsilon>0$ wybieramy odpowiadające mu $N$.
Jeśli $i>j\ge N$, to
$\|S_i-S_j\|=\|\sum_{n=j+1}^i x_n\|<\varepsilon$,
bo $i\ge j+1\ge N$. Dla $i<j$ korzystamy z symetrii
normy różnicy, a dla $i=j$ różnica jest zerowa.
Ciąg $(S_n)$ jest zatem ciągiem Cauchy'ego w $E$.
Jeśli $E$ jest zupełna, istnieje $S\in E$ z $S_n\to S$,
co jest definicją zbieżności szeregu.
\end{proof}

\begin{przyklad}[Wyrazy dążące do zera nie wystarczają]
\label{ma2:ex:vector-harmonic}
Jeżeli $v\in E$ i $v\ne0$, szereg $\sum_{n\ge1}v/n$
jest rozbieżny, chociaż $\|v/n\|=\|v\|/n\to0$.
Istotnie dla każdego $N\ge1$ mamy
\[
 \left\|\sum_{n=N+1}^{2N}\frac vn\right\|
 =\|v\|\sum_{n=N+1}^{2N}\frac1n\ge\frac{\|v\|}{2}.
\]
Wybierając $\varepsilon=\|v\|/2>0$, widzimy, że dla
każdego proponowanego indeksu w kryterium Cauchy'ego istnieje
ogon skończony o normie co najmniej $\varepsilon$.
Z twierdzenia \ref{ma2:thm:series-cauchy} szereg jest rozbieżny.
\end{przyklad}

\section{Rachunek sum i operatory ciągłe}
\begin{twierdzenie}[Pominięcie skończonego początku i liniowość]
\label{ma2:thm:series-linear-tail}
Dla każdego całkowitego $m\ge1$ szereg $\sum_{n\ge1}x_n$
jest zbieżny wtedy i tylko wtedy, gdy zbieżny jest
$\sum_{n\ge m}x_n$. W przypadku zbieżności
\[
 \sum_{n=1}^\infty x_n=\sum_{n=1}^{m-1}x_n+
                       \sum_{n=m}^\infty x_n.
\]
Jeśli szeregi $\sum x_n$ i $\sum y_n$ w $E$ są zbieżne
oraz $\alpha,\beta\in\K$, to szereg $\sum(\alpha x_n+\beta y_n)$
jest zbieżny i
\[
 \sum_{n=1}^\infty(\alpha x_n+\beta y_n)
 =\alpha\sum_{n=1}^\infty x_n+\beta\sum_{n=1}^\infty y_n.
\]
\end{twierdzenie}
\begin{proof}
Dla $N\ge m$ suma od $m$ do $N$ jest równa $S_N-S_{m-1}$.
Jeśli $S_N\to S$, to
$\|(S_N-S_{m-1})-(S-S_{m-1})\|=\|S_N-S\|\to0$.
Odwrotnie, jeśli sumy od $m$ do $N$ dążą do $A$, to
$\|S_N-(A+S_{m-1})\|=\|(S_N-S_{m-1})-A\|\to0$.
Daje to obie implikacje i wzór, również dla $m=1$,
gdy $S_0=0$.
Niech $S_N\to S$ i $T_N=\sum_{n=1}^N y_n\to T$.
Suma częściowa kombinacji liniowej wynosi
$\alpha S_N+\beta T_N$, a jej odległość od
$\alpha S+\beta T$ jest nie większa niż
$|\alpha|\|S_N-S\|+|\beta|\|T_N-T\|\to0$.
Nie potrzeba tu zupełności $E$, ponieważ obie granice
istnieją z założenia.
\end{proof}

\begin{twierdzenie}[Przenoszenie sumy przez operator]
\label{ma2:thm:series-operator}
Niech $F$ będzie przestrzenią unormowaną nad $\K$,
a $T:E\to F$ ciągłym operatorem liniowym.
Jeśli $\sum x_n=S$ w $E$, to $\sum Tx_n=TS$ w $F$.
\end{twierdzenie}
\begin{proof}
Z liniowości $\sum_{n=1}^N Tx_n=TS_N$.
Z twierdzenia \ref{prop:operator-norm12} mamy
$\|TS_N-TS\|_F\le\|T\|\|S_N-S\|_E\to0$.
Granica $TS$ należy do $F$ niezależnie od zupełności $F$.
\end{proof}

\begin{twierdzenie}[Współrzędne w skończonym wymiarze]
\label{ma2:thm:series-coordinates}
Niech $E$ będzie skończenie wymiarowa, z bazą
$(e_j)_{j=1}^d$, gdzie $d\ge0$, i niech
$x_n=\sum_{j=1}^d a_{nj}e_j$, $a_{nj}\in\K$.
Szereg $\sum_n x_n$ jest zbieżny wtedy i tylko wtedy,
gdy każdy z $d$ szeregów skalarnych $\sum_n a_{nj}$
jest zbieżny. Wtedy
\[
 \sum_{n=1}^\infty x_n
 =\sum_{j=1}^d\left(\sum_{n=1}^\infty a_{nj}\right)e_j.
\]
\end{twierdzenie}
\begin{proof}
Dla $d=0$ każdy wyraz i każda suma są zerowe, a warunek
zbieżności szeregów współrzędnych jest spełniony w sposób pusty.
Załóżmy więc $d\ge1$.
Z dowodu twierdzenia \ref{ma2:thm:finite} istnieją $c,C>0$
takie, że dla każdego $z\in\K^d$ zachodzi
$c\max_j|z_j|\le\|\sum_j z_je_j\|\le C\max_j|z_j|$.
Oznaczmy $s_{Nj}=\sum_{n=1}^N a_{nj}$.
Z rozdzielności skończonych sum $S_N=\sum_j s_{Nj}e_j$.
Jeśli $S_N\to S=\sum_j s_je_j$, dolne oszacowanie daje
$|s_{Nj}-s_j|\le\|S_N-S\|/c\to0$ dla każdego $j$.
Odwrotnie, jeśli $s_{Nj}\to s_j$ dla wszystkich $j$, to
dla $\varepsilon>0$ dobieramy dla każdego $j$ indeks $N_j$
z $|s_{Nj}-s_j|<\varepsilon/C$ dla $N\ge N_j$.
Dla $N\ge\max_j N_j$ górne oszacowanie daje
$\|S_N-\sum_j s_je_j\|<\varepsilon$.
To dowodzi zbieżności i wzoru.
\end{proof}

\section{Szeregi geometryczne, teleskopowe i grupowanie}
\begin{twierdzenie}[Szereg geometryczny pomnożony przez wektor]
\label{ma2:thm:vector-geometric}
Niech $v\in E$ i $q\in\K$, przy czym $q^0=1$ również dla $q=0$.
Jeżeli $|q|<1$, to
\[
 \sum_{n=0}^{\infty}q^n v=\frac{1}{1-q}v.
\]
Jeśli $v\ne0$, warunek $|q|<1$ jest także konieczny.
Dla $v=0$ szereg jest zbieżny dla każdego $q$.
\end{twierdzenie}
\begin{proof}
Dla $|q|<1$ mamy $q\ne1$; tożsamość skończonego szeregu
geometrycznego daje
$\sum_{n=0}^N q^n v=(1-q^{N+1})v/(1-q)$.
Odległość od $v/(1-q)$ wynosi
$|q|^{N+1}\|v\|/|1-q|$, co dąży do zera, ponieważ
$0\le |q|<1$. Nie używamy zupełności $E$.
Jeśli $v\ne0$ i $|q|\ge1$, normy wyrazów
$|q|^n\|v\|$ są nie mniejsze niż $\|v\|>0$.
Warunek konieczny z twierdzenia \ref{ma2:thm:series-cauchy}
(po pominięciu wyrazu o indeksie $0$) nie jest spełniony.
Jeśli $v=0$, wszystkie sumy częściowe są zerowe.
\end{proof}

\begin{twierdzenie}[Szereg teleskopowy]
\label{ma2:thm:vector-telescoping}
Dla ciągu $(u_n)$ w $E$ szereg $\sum_{n\ge1}(u_n-u_{n+1})$
jest zbieżny wtedy i tylko wtedy, gdy $(u_n)$ jest zbieżny.
Jeśli $u_n\to u$, suma wynosi $u_1-u$.
\end{twierdzenie}
\begin{proof}
Dla $N\ge1$ sumy pośrednich wyrazów znoszą się, więc
$\sum_{n=1}^N(u_n-u_{n+1})=u_1-u_{N+1}$.
Jeśli $u_n\to u$, norma różnicy tej sumy od $u_1-u$
wynosi $\|u_{N+1}-u\|\to0$.
Jeśli natomiast suma częściowa dąży do $S$, to
$\|u_{N+1}-(u_1-S)\|=\|S-(u_1-u_{N+1})\|\to0$.
Pominięcie pierwszego wyrazu nie zmienia granicy ciągu,
zatem $u_n\to u_1-S$.
\end{proof}

\begin{twierdzenie}[Grupowanie kolejnych wyrazów]
\label{ma2:thm:series-grouping}
Niech $0=k_0<k_1<k_2<\cdots$ będą liczbami całkowitymi
i niech $b_j=\sum_{n=k_{j-1}+1}^{k_j}x_n$ dla $j\ge1$.
Jeżeli $\sum x_n=S$, to $\sum b_j=S$.
Odwrotnie, jeśli $\sum b_j=S$ oraz
\[
 \max_{k_{j-1}\le m\le k_j}
 \left\|\sum_{n=k_{j-1}+1}^{m}x_n\right\|\longrightarrow0,
\]
to $\sum x_n=S$.
\end{twierdzenie}
\begin{proof}
Z rozłączności kolejnych przedziałów indeksów
$\sum_{j=1}^J b_j=\sum_{n=1}^{k_J}x_n=S_{k_J}$.
Ponieważ $k_J\ge J$, jest to podciąg sum częściowych,
a z $S_N\to S$ wynika $S_{k_J}\to S$.
Dla implikacji odwrotnej oznaczmy maksimum z tezy przez $d_j$.
Dla każdego $N\ge1$ istnieje dokładnie jeden $j$ z
$k_{j-1}<N\le k_j$, ponieważ $(k_j)$ jest ściśle rosnący
i nieograniczony. Wtedy
\[
 \|S_N-S\|\le\|S_{k_{j-1}}-S\|+
              \|S_N-S_{k_{j-1}}\|
 \le\|S_{k_{j-1}}-S\|+d_j.
\]
Dla $\varepsilon>0$ wybieramy $J$ tak duże, że dla $j\ge J$
obie wielkości po prawej stronie są mniejsze niż
$\varepsilon/2$ (pierwsza z założenia zbieżności szeregu
zgrupowanego, druga z $d_j\to0$).
Dla $N>k_{J-1}$ odpowiedni indeks $j$ jest co najmniej $J$,
więc $\|S_N-S\|<\varepsilon$.
\end{proof}
\begin{przyklad}[Grupowanie nie musi zachowywać rozbieżności]
\label{ma2:ex:grouping-divergence}
Dla $v\ne0$ połóżmy $x_{2j-1}=v$, $x_{2j}=-v$.
Sumy częściowe o indeksach parzystych wynoszą $0$,
a o nieparzystych $v$. Ponieważ $v\ne0$, oba podciągi
mają różne granice, więc szereg nie jest zbieżny.
Przy $k_j=2j$ każdy wyraz $b_j$ jest zerowy,
a szereg zgrupowany jest zbieżny do $0$.
Warunek maksimum w twierdzeniu \ref{ma2:thm:series-grouping}
nie zachodzi, ponieważ w każdym bloku maksimum wynosi $\|v\|$.
\end{przyklad}

\section{Zbieżność bezwzględna i oszacowania reszty}
\begin{definicja}[Zbieżność bezwzględna]
\label{ma2:def:absolute-series}
Szereg $\sum x_n$ nazywamy bezwzględnie zbieżnym, jeżeli
zbieżny jest szereg liczbowy $\sum\|x_n\|$.
Sama definicja nie wymaga zupełności $E$.
\end{definicja}
\begin{twierdzenie}
\label{ma2:thm:absolute-series}
Jeżeli $E$ jest przestrzenią Banacha i $\sum\|x_n\|<\infty$,
to $\sum x_n$ jest zbieżny. Dla jego sumy $S$ i każdego
$N\ge0$ zachodzi oszacowanie
\[
 \|S-S_N\|\le\sum_{n=N+1}^\infty\|x_n\|.
\]
\end{twierdzenie}
\begin{proof}
Dla $q\ge p$ nierówność
$\|\sum_{n=p}^q x_n\|\le\sum_{n=p}^q\|x_n\|$
wynika przez indukcję: dla jednego składnika mamy równość,
a dołączenie następnego wyrazu zwiększa prawą stronę o jego normę
na mocy nierówności trójkąta.
Dla $\varepsilon>0$ kryterium Cauchy'ego dla zbieżnego
szeregu liczbowego $\sum\|x_n\|$ daje $M$, dla którego
$\sum_{n=p}^q\|x_n\|<\varepsilon$ przy $q\ge p\ge M$.
Stąd spełniony jest warunek z twierdzenia
\ref{ma2:thm:series-cauchy}. Zupełność $E$ daje istnienie $S$.
Dla $L>N$ mamy
$\|S_L-S_N\|\le\sum_{n=N+1}^L\|x_n\|$.
Przy ustalonym $N$ lewa strona dąży do $\|S-S_N\|$
z ciągłości normy (twierdzenie \ref{ma2:thm:metric}),
a prawa do sumy ogona zbieżnego szeregu liczbowego.
Przechodząc do granicy w nierówności nieostrej, dostajemy tezę.
\end{proof}

\begin{twierdzenie}[Charakteryzacja zupełności]
\label{ma2:thm:series-completeness}
Przestrzeń unormowana $E$ jest przestrzenią Banacha wtedy
i tylko wtedy, gdy każdy szereg bezwzględnie zbieżny w $E$
jest zbieżny w $E$.
\end{twierdzenie}
\begin{proof}
Pierwsza implikacja wynika z twierdzenia \ref{ma2:thm:absolute-series}.
Załóżmy własność zbieżności wszystkich takich szeregów
i niech $(y_n)$ będzie dowolnym ciągiem Cauchy'ego w $E$.
Dla każdego $j\ge1$ istnieje $M_j\ge1$ takie, że
$\|y_n-y_m\|<2^{-j}$ dla $n,m\ge M_j$.
Możemy wybrać $M_j$ jako najmniejszy dodatni indeks
o podanej własności. Wybieramy $n_1=M_1$ i rekurencyjnie
$n_{j+1}=\max\{n_j+1,M_{j+1},M_j\}$.
Z definicji $n_{j+1}>n_j$ i $n_{j+1}\ge M_{j+1}$.
Ponieważ $n_1=M_1$, indukcyjnie $n_j\ge M_j$ dla każdego $j$.
Ponadto $n_{j+1}\ge M_j$, zatem $n_j,n_{j+1}\ge M_j$,
więc dla $h_j=y_{n_{j+1}}-y_{n_j}$ mamy $\|h_j\|<2^{-j}$.
Szereg $\sum\|h_j\|$ jest zbieżny przez porównanie
z szeregiem geometrycznym. Z założenia istnieje
$h=\sum_{j=1}^\infty h_j$ w $E$.
Suma teleskopowa daje $y_{n_{J+1}}=y_{n_1}+\sum_{j=1}^J h_j$,
a zatem $y_{n_{J+1}}\to y=y_{n_1}+h$.
Ustalmy $\varepsilon>0$. Wybieramy $M$ z
$\|y_n-y_m\|<\varepsilon/2$ dla $n,m\ge M$.
Ponieważ $n_{J+1}\to\infty$ i $y_{n_{J+1}}\to y$,
wybieramy $J$ z $n_{J+1}\ge M$ oraz
$\|y_{n_{J+1}}-y\|<\varepsilon/2$.
Dla każdego $n\ge M$ dostajemy
$\|y_n-y\|\le\|y_n-y_{n_{J+1}}\|+\|y_{n_{J+1}}-y\|<\varepsilon$.
Dowolny ciąg Cauchy'ego ma więc granicę w $E$.
\end{proof}

\begin{przyklad}[Konieczność zupełności]
\label{ma2:ex:incomplete-series}
W przestrzeni $c_{00}$ z normą supremum niech $e_n$
oznacza ciąg mający $1$ na $n$-tej współrzędnej i zera na
pozostałych. Wtedy $x_n=2^{-n}e_n$ należy do $c_{00}$ i
$\sum\|x_n\|_\infty=\sum 2^{-n}=1$.
Sumy częściowe $S_N$ mają w $\ell^\infty$ granicę
$x=(2^{-j})_{j\ge1}$, ponieważ
$\|S_N-x\|_\infty=\sup_{j>N}2^{-j}=2^{-(N+1)}\to0$.
Gdyby szereg był zbieżny w $c_{00}$, granica byłaby
również jego granicą w $\ell^\infty$, gdzie norma jest ta sama.
Jednoznaczność dawałaby $x\in c_{00}$, co jest niemożliwe,
ponieważ $x_j>0$ dla każdego $j$.
\end{przyklad}

\section{Kryteria porównawcze, ilorazowe i pierwiastkowe}
\begin{wniosek}[Majoranta liczbowa]
\label{ma2:cor:series-majorant}
Niech $E$ będzie przestrzenią Banacha. Jeśli $a_n\ge0$,
$\sum a_n<\infty$ i od pewnego indeksu $\|x_n\|\le a_n$,
to $\sum x_n$ jest bezwzględnie zbieżny i zbieżny.
Jeśli dla całkowitego $N\ge0$ nierówność zachodzi dla wszystkich $n>N$, to
$\|S-S_N\|\le\sum_{n>N}a_n$.
\end{wniosek}
\begin{proof}
Wybierzmy $m\ge1$ tak, aby $\|x_n\|\le a_n$ dla $n\ge m$.
Dla $L\ge m$ mamy
\[
 \sum_{n=1}^L\|x_n\|\le
 \sum_{n=1}^{m-1}\|x_n\|+\sum_{n=m}^L a_n
 \le \sum_{n=1}^{m-1}\|x_n\|+\sum_{n=1}^\infty a_n.
\]
To samo skończone ograniczenie działa dla $L<m$,
bo normy są nieujemne. Sumy częściowe szeregu norm są
niemalejące i ograniczone z góry, więc są zbieżne
na mocy twierdzenia o zbieżności ciągu monotonicznego
z tomu I. Z twierdzenia \ref{ma2:thm:absolute-series}
istnieje suma $S$ szeregu wektorowego.
Jeżeli $\|x_n\|\le a_n$ dla wszystkich $n>N$, to dla $L>N$
mamy $\sum_{n=N+1}^L\|x_n\|\le\sum_{n=N+1}^L a_n$.
Obie strony mają skończone granice przy $L\to\infty$,
więc oszacowanie reszty z twierdzenia
\ref{ma2:thm:absolute-series} daje
$\|S-S_N\|\le\sum_{n=N+1}^\infty\|x_n\|
\le\sum_{n=N+1}^\infty a_n$.
\end{proof}

\begin{twierdzenie}[Kryteria d'Alemberta i Cauchy'ego]
\label{ma2:thm:vector-ratio-root}
Niech $E$ będzie przestrzenią Banacha.
\begin{enumerate}[label=(\roman*)]
\item Jeżeli od pewnego indeksu $x_n\ne0$, to dla
$r_n=\|x_{n+1}\|/\|x_n\|$ warunek $\limsup r_n<1$
daje zbieżność bezwzględną i zbieżność szeregu.
Warunek $\liminf r_n>1$ daje rozbieżność.
\item Dla $L=\limsup_{n\to\infty}\|x_n\|^{1/n}$
warunek $L<1$ daje zbieżność bezwzględną i zbieżność,
a $L>1$ daje rozbieżność. W tym przypadku dopuszczamy
wyrazy zerowe.
\end{enumerate}
\end{twierdzenie}
\begin{proof}
Granice dolną i górną rozumiemy w rozszerzonej prostej
rzeczywistej, zgodnie z tomem I (rozdział 08). Korzystamy
z kryterium progowego: jeżeli $\limsup b_n<q$, to od
pewnego indeksu $b_n<q$; jeżeli $q<\liminf b_n$, to od
pewnego indeksu $q<b_n$; jeżeli $q<\limsup b_n$, to
dla każdego $M$ istnieje $n\ge M$ z $q<b_n$.
W części (i), jeśli $\limsup r_n<1$, wybieramy
$0<q<1$ z $\limsup r_n<q$. Od pewnego $m$ mamy
$\|x_{n+1}\|\le q\|x_n\|$ dla $n\ge m$.
Indukcyjnie $\|x_{m+k}\|\le q^k\|x_m\|$ dla $k\ge0$:
dla $k=0$ mamy równość, a przejście do $k+1$ wynika
z powyższego oszacowania kolejnych norm.
Majoranta geometryczna i wniosek \ref{ma2:cor:series-majorant}
dają zbieżność.
Jeśli $\liminf r_n>1$, wybieramy skończone $q>1$ mniejsze
od granicy dolnej (również gdy jest nieskończona).
Od pewnego $m$ mamy $r_n>q$ i $x_n\ne0$.
Dla $k=0$ mamy równość $\|x_{m+k}\|=\|x_m\|$.
Jeżeli $\|x_{m+k}\|\ge q^k\|x_m\|$, to
$\|x_{m+k+1}\|>q\|x_{m+k}\|\ge q^{k+1}\|x_m\|$.
Zatem indukcyjnie $\|x_{m+k}\|\ge q^k\|x_m\|\ge\|x_m\|>0$
dla każdego $k\ge0$. Z twierdzenia \ref{ma2:thm:series-cauchy}
szereg jest rozbieżny.
W części (ii), jeśli $L<1$, wybieramy $0<q<1$ z $L<q$.
Od pewnego $m$ mamy $\|x_n\|^{1/n}<q$, czyli
$\|x_n\|<q^n$ dla $n\ge m$ (potęgowanie stopnia $n$ zachowuje
porządek na liczbach nieujemnych).
Stosujemy wniosek \ref{ma2:cor:series-majorant}.
Jeśli $L>1$, wybieramy skończone $q>1$ z $q<L$.
Dla każdego $M$ istnieje $n\ge M$ z $\|x_n\|^{1/n}>q$,
stąd $\|x_n\|>q^n\ge1$. Wyrazy nie dążą do zera,
więc szereg jest rozbieżny.
\end{proof}
\begin{przyklad}[Kontrola błędu obcięcia]
\label{ma2:ex:series-geometric-bound}
Jeśli $C\ge0$, $0\le q<1$ i $\|x_n\|\le Cq^n$ dla $n\ge1$,
to w przestrzeni Banacha z wniosku
\ref{ma2:cor:series-majorant} istnieje suma $S$ i
\[
 \|S-S_N\|\le C\sum_{n=N+1}^\infty q^n
              =\frac{Cq^{N+1}}{1-q}\qquad(N\ge0).
\]
Dla $q=0$ wszystkie $x_n$ są zerowe, więc oszacowanie
nadal zachodzi. Dla $C=1$, $q=1/2$ i $N=10$ błąd jest
nie większy niż $2^{-10}<1/1000$.
\end{przyklad}

\begin{uwaga}
Rozbieżność wykazana przez warunek konieczny nie wymaga
zupełności. W tym rozdziale nie utożsamiamy zbieżności
bezwzględnej ze zbieżnością bezwarunkową i nie zmieniamy
kolejności wyrazów. Przestawienia, szeregi warunkowo zbieżne
i szeregi bezwarunkowo zbieżne będą tematem rozdziału 03.
\end{uwaga}
